\subsection{معماری و پیاده‌سازی}
کلاس اصلی پیاده‌سازی‌شده در چارچوب پروژه تحت عنوان \lr{GraniteHybridBlock} در فایل \lr{Codes/framework/attentions/granite\_hybrid.py} قرار داشته و از کلاس مرجع \lr{BaseAttention} ارث‌بری می‌کند. این ماژول همچنین شامل دو زیرماژول کلیدی \lr{Mamba2} و \lr{MixtureOfExperts} است.

مشخصات کلیدی ساختار پیاده‌سازی عبارتند از:
\begin{itemize}
    \item \textbf{آرگومان‌های کانفیگ و سازنده:} سازنده کلاس آرگومان‌های \lr{config}، \lr{rank}، \lr{dim}، \lr{num\_heads}، \lr{num\_experts=4}، \lr{top\_k=2}، \lr{ratio=2}، \lr{lora\_rank=16} و \lr{lora\_alpha=32.0} را جهت پیکربندی دریافت می‌کند.
    \item \textbf{زیرماژول Mamba2:} شامل کانولوشن ۱ بعدی با طول پنجره زمانی ۴ (\lr{kernel\_size=4})، پروجکشن گام زمانی دلتا \lr{dt\_proj} و پارامترهای پایداری جهت ممانعت از سرریز گرادیان در توابع نمایی (\lr{torch.clamp}).
    \item \textbf{زیرماژول MixtureOfExperts:} شبکه‌ای با ۴ خبره خطی مستقل دولایه مجهز به فعال‌ساز \lr{SiLU} و مسیریاب گیت \lr{router} که با عملگر \lr{torch.topk} دو خبره برتر را انتخاب می‌کند.
    \item \textbf{تزریق و انطباق چندمعماری:} توابع کمکی \lr{copy\_weights} و \lr{inject\_granite\_hybrid} سازگاری کامل لایه را با ساختارهای مختلف انکودر و دیکودر از قبیل \lr{DistilBERT}، \lr{RoBERTa}، \lr{GPT-2} و \lr{Qwen} محقق می‌سازند.
\end{itemize}

\subsection{محدودیت‌ها و تحلیل شکست}
با وجود ظرفیت نظری بسیار بالای معماری، تحلیل جزئیات حاکی از نکاتی در استقرار است:
\begin{enumerate}
    \item \textbf{افت دقت در \lr{RoBERTa} و \lr{GPT-2} (\lr{PASS\_LOW}):} دقت مدل در \lr{RoBERTa} برابر با ۵۳.۳۳٪ و در \lr{GPT-2} برابر با ۵۱.۷۲٪ است. علت این امر به ناهمگونی در توزیع نمایش فضایی شاخه \lr{Mamba-2} و شاخه توجه بازمی‌گردد. در فاز انطباق سبک بدون فاین‌تیونینگ تمام پارامترها، وزن‌های پیش‌آموزش‌دیده توجه متراکم با حالات تصادفی اولیه کانولوشن و ماتریس انتقال حالت $\bar{\mathbf{A}}$ تداخل مخرب پیدا کرده و منجر به سوگیری تخمین به سمت توزیع یکنواخت می‌گردد.
    \item \textbf{افزایش ردپای حافظه سخت‌افزاری:} در مدل‌های بزرگ‌تر، نگهداری هم‌زمان پارامترهای خبره‌ها در حافظه \lr{VRAM} (به ویژه ۶۵۹۱ مگابایت در \lr{Qwen}) نشان می‌دهد که علی‌رغم کاهش مرتبه نظری توجه، فعال‌سازی‌های خبره‌های \lr{MoE} و نگهداری تانسورهای کانولوشن زمانی باعث ایجاد گلوگاه حافظه در کارت‌های گرافیک با ظرفیت متوسط می‌شود.
\end{enumerate}
